\chapter{在线教学}\label{s:online}

\begin{quote}

  如果你用机器人来教书，你就是在把人教成机器人。\\
  —— 出处不一

\end{quote}

技术已经多次改变教与学。
在 1800 年代初黑板被引入学校之前，
老师没法把一个即兴的例子、
一张图、
或一道练习一次分享给全班。
便宜、
可靠、
易用、
灵活的黑板，让老师能以很快的速度、很大的规模做以前只能慢慢做、零零碎碎做的事。
同样，
手持摄像机给田径训练带来了革命，
就像十年前录音机给音乐教学带来的革命一样。

许多把互联网推进课堂的人并不了解这段历史，
也没有意识到他们的尝试，只是
\hreffoot{http://teachingmachin.es/timeline.html}{一长串用机器来教学的尝试}中最新的一次~\cite{Watt2014}。
从印刷机到广播、电视，
再到台式电脑和移动设备，
每一种分享知识的新方式，都催生了一波激进的乐观主义者，
他们相信教育已经坏掉、而技术能把它修好。
然而，
教育科技叫得最响的倡导者，往往对「教育」的了解还不如对「科技」，
而在他们的辞藻背后，
很多人更多是被盈利前景驱动，
而不是被赋能学习者的愿望驱动。

今天的争论常常把「在线」和「自动化」混为一谈，从而变得浑浊。
运行得好，
十几个人在视频聊天里一起解决一个问题，
感觉就和任何别的小组讨论一样。
反过来，
一队助教照着死板的评分标准批改几百份卷子，
那还不如说是一堆 Perl 脚本。\index{Perl（贬义指代）}
所以本章先看使用录制视频和自动评分练习的完全自动化在线教学，
然后再探讨一些替代性的混合模式。

\seclbl{MOOC}{s:online-moocs}

用互联网重塑教育最引人注目的努力，
是\gref{g:mooc}{大规模开放在线课程}，即 MOOC。
这个词由 David Cormier 在 2008 年创造，\index{Cormier, David}
用来描述 George Siemens\index{Siemens, George}
和 Stephen Downes\index{Downes, Stephen}组织的一门课程。
那门课程基于一种\grefdex{g:connectivism}{联通主义}{联通主义}的学习观，
认为知识是分布式的，
学习就是发现、建立和修剪连接的过程。

「MOOC」这个词很快被一些课程的创造者据为己有，
那些课程更像是传统课堂的「中心—辐射」模式，
老师处在中心、定义目标，
学习者则被视为知识的接受者或复制者。
使用最初联通主义模式的课，现在有时被称为「cMOOC」，
而把控制集中起来的课则被称为「xMOOC」。
（后者有时也被叫做「MESS」，
即「大规模强化的台上圣人」。）

五年前，
你走过任何一所主要大学校园，
都能听到有人在谈论 MOOC 会如何给教育带来革命、
会如何毁掉教育、
或者两者兼有。
MOOC 会给学习者提供更广泛的课程，
让他们在自己方便的时候学习，
而不必让学习去迁就别人的日程。

但 MOOC 的成效远不如它更热情的拥趸所预测的那样~\cite{Ubel2017}。
一个原因是，
录制内容对许多新手效果不佳，
因为它无法清除他们各自的误解（\chapref{s:models}）：
如果第一次没听懂某个讲解，
通常也没有另一个版本可选。
另一个原因是，
要让 MOOC 配得上「大规模」所需的自动评分，
在布鲁姆分类法里只在最低的那些层次上才管用（\secref{s:process-objectives}）。\index{布鲁姆分类法}
现在也已经很清楚：
在 MOOC 里，学习者得更多地自己扛起保持专注的责任、
在线学习的非人际性会助长不文明行为、打击人，
而且「对所有人开放」实际上意味着
「对所有足够富裕、有高速网络和大量空闲时间的人开放」。

\cite{Marg2015}考察了 76 门不同主题的 MOOC，
发现虽然材料的组织和呈现不错，
但课程设计的质量很差。
更贴近我们的是，
\cite{Kim2017}研究了三十个流行的在线编程教程，
发现它们大体上以相同的方式教相同的内容：
自下而上，
从低层次的编程概念讲起，逐步上升到高层次目标。
大多数要求学习者写程序、并提供某种形式的即时反馈，
但这些反馈通常非常浅薄。
很少有教程解释概念在何时、为何有用
（也就是没有展示如何迁移知识），
也没有为常见错误提供指引；
除了根据年龄做的粗浅区分之外，
没有一个教程会根据先前的编程经验或学习者目标来个性化课程。

\begin{aside}{个性化学习}
  很少有哪个术语像\gref{g:personalized-learning}{个性化学习}这样被以五花八门的方式使用和滥用。
  对大多数教育科技倡导者来说，
  它意味着根据学习者表现动态调整课程进度，
  这样如果某人连续答对了几道题，
  计算机就会跳过后面的一些题。

  这么做能带来
  \hreffoot{https://www.rand.org/pubs/research\_briefs/RB9994.html}{有限的改善}，
  但本可以做得更好。
  例如，
  如果许多学习者觉得某个主题难，
  老师可以为那一点准备多种不同的讲解方式，
  而不是在单一路径上加速。
  这样，
  如果一种讲法没有引起共鸣，
  还有别的可选。
  然而，
  这要求老师做多得多的设计工作，
  这也许就是它没能流行起来的原因。
  而且即便它有效，
  效果也很可能远小于它某些倡导者所相信的。
  在小学，一位好老师能让学年末的表现差出 0.1—0.15 个标准差~\cite{Chet2014}
  （简要总结见\hreffoot{http://educationnext.org/in-schools-teacher-quality-matters-most-coleman/}{这篇文章}）。
  认为任何形式的自动化短期内能超过这一点，是不现实的。
\end{aside}

那么，互联网\emph{应该}怎样用于教与学技术技能呢？
它的利与弊是：\index{在线学习!利与弊}

\begin{description}

\item[学习者能比以往更快地获取更多课程。]
  当然，前提是
  搜索引擎认为这些课程值得索引、
  他们的网络服务提供商和政府没有屏蔽它、
  而且真相没有被一片耗尽注意力的虚假信息淹没。

\item[学习者能比以往获取\emph{更好}的课程，]
  除非他们被引向二流材料，
  以便把财富从「没有的人」重新分配给「有的人」~\cite{McMi2017}。
  同样值得记住的是，稀缺会抬高感知价值，
  所以随着在线教育变便宜，
  它会越来越变成「人人都希望给别人的孩子用」的东西。

\item[学习者也能接触到比以往多得多的人。]
  但前提是这些学习者确实有使用所需技术的条件、
  用得起它、
  并且不会因为骚扰而被赶下线、
  或因为不符合说话最响的那群人的社会规范而被边缘化。
  实际上，
  大多数 MOOC 用户来自有保障、富裕的背景~\cite{Hans2015}。

\item[老师能获得关于学习者如何学习的细致得多的洞察。]
  只要学习者做的事适合大规模自动化分析，
  并且要么不反对课堂里的监视，
  要么其反对没有足够的分量让人在乎。

\end{description}

\cite{Marg2015,Mill2016a,Nils2017}描述了如何在避免上述弊端的同时，放大其中的好处：

\begin{description}

\item[把截止日期设得频繁并广为告知，]
  并严格执行，好让学习者进入一个学习节奏。

\item[把现场讲座这类全班同步活动降到最少，]
  免得人们因为日程冲突而错过东西。

\item[让学习者对集体知识做出贡献，]
  例如一起记笔记（\secref{s:classroom-notetaking}）、
  担任课堂记录员、
  或向共享题库贡献题目（\secref{s:individual-peer}）。

\item[鼓励或要求学习者以小组完成一部分工作，]
  这些小组\emph{确实}有同步在线活动，
  比如每周一次在线讨论。
  这有助于学习者保持投入和动力，又不会制造太多日程上的麻烦。
  （让这些讨论公平而有成效的一些技巧，见\appref{s:meetings}。）

\item[制定、公布并执行一份行为准则，]
  好让每个人都能真正参与在线讨论（\secref{s:classroom-coc}）。

\item[用大量简短的课程片段，而不是几段讲座长度的长块，]
  以尽量减少认知负荷、
  并为形成性评价提供大量机会。
  这也有助于维护：
  如果你的视频都很短，
  需要维护时直接重录就行，
  这往往比给较长的视频打补丁更便宜。

\item[用视频来吸引，而不是用来说教。]
  撇开无障碍问题不谈（\secref{s:motivation-accessibility}），
  学习者读得比你讲得快。
  这条规则的一个例外是，
  视频其实是教「动词」（动作）的最佳方式：
  短小的录屏能教人怎么用编辑器、
  在调试器里单步走代码，
  等等，比带文字的截图更有效。

\item[尽早识别并清除误解。]
  如果数据显示学习者在课程的某些部分卡住了，
  就为那些点制作替代讲解，
  并给他们额外的练习去操练。

\end{description}

这一切总得用某种方式落实，\index{在线学习!落实}
这意味着你需要某种教学平台。
你既可以用一个一体化的\gref{g:lms}{学习管理系统}，
比如 \hreffoot{http://moodle.org}{Moodle} 或 \hreffoot{https://www.sakaiproject.org/}{Sakai}，
也可以自己拼装一套，
用 \hreffoot{http://slack.com}{Slack} 或 \hreffoot{https://zulipchat.com/}{Zulip} 做聊天，
用 \hreffoot{http://hangouts.google.com}{Google Hangouts}
或 \hreffoot{https://appear.in/}{appear.in} 做视频对话，
用 \hreffoot{https://wordpress.org/}{WordPress}、
\hreffoot{http://docs.google.com}{Google Docs}、
或随便哪种 wiki 系统做协作写作。
如果你刚开始，
就挑最容易搭建和管理、
而且你的学习者最熟悉的那种。
如果二者只能选一，
第二个考虑比第一个更重要：
你指望人们在你课上学很多东西，
那么你学会用他们最顺手的工具，才算公平。

搭一个学习平台是必要的，但还不够：
如果你想让学习者茁壮成长，
你需要创建一个社区。
有几百本书和演讲讲怎么做到这一点，
但大多数是基于作者的亲身经验。
\cite{Krau2016}是个令人欣慰的例外：
尽管它成书于 Twitter 和 Facebook 加速堕落成武器化的滥用和虚假信息之前，
它的大部分发现仍然有用。
\cite{Foge2005}也充满了关于学习者可能希望加入的实践共同体的实用技巧；
我们在\chapref{s:community}里探讨它的一些想法。

\begin{aside}{「做」的自由与「免于」的自由}
  Isaiah Berlin 1958 年的文章\index{Berlin, Isaiah}
  「\hreffoot{https://en.wikipedia.org/wiki/Two\_Concepts\_of\_Liberty}{两种自由概念}」
  区分了积极自由，
  即真正能做某件事的能力，
  和消极自由，
  即没有规则说你不许做。
  在线讨论通常提供的是消极自由
  （没人阻止你说出你的想法）
  而不是积极自由
  （许多人其实没法被听到）。
  应对这一点的一种办法是引入某种限流，
  比如每天每个讨论串里每位学习者只允许发一条消息。
  这样就让有话要说的人有机会说出来，
  同时为别人腾出说话的空间。
\end{aside}

人们对在线教学还有另一个担忧，是作弊。
日常的不诚实，在线课上并不比面对面更多~\cite{Beck2014}，
但请别人代写期末考试的诱惑、
以及核查这件事是否发生的难度，
是教育机构一直不愿给纯在线课程学分的原因之一。
远程监考是可能的，
但在为此投入之前，
先读读~\cite{Lang2013}：
它探讨了学习者为什么、以及怎样作弊，
以及课程可以怎样设计，免得给他们作弊的理由。

\seclbl{视频}{s:online-video}

大多数 MOOC 的一个突出特点，是使用录制的视频讲座。
它们可以是有效的：
正如\chapref{s:performance}里提到的，
一种叫直接教学、\index{直接教学}
以精确呈现精心设计的脚本为基础的教学技巧，已被反复证明是有效的~\cite{Stoc2018}。
然而，
直接教学的脚本必须非常仔细地设计、
测试、
和打磨，
这是许多 MOOC 不愿、也无法做出的投入。
给网页或幻灯片做个小改动只需几分钟；
给一个短视频哪怕做个小改动也要一小时以上，
所以老师根据反馈采取行动的成本可能承担不起。
而且即便做得很好，
视频也必须和活动结合起来才有益：
\cite{Koed2015}估计，
「{\ldots}多做带来的学习收益{\ldots}
是多看或多读的六倍以上。」

如果你教编程，
可以用录屏代替幻灯片，\index{录屏}
因为它们提供了一些和现场编程相同的好处（\secref{s:performance-live}）。
\cite{Chen2009}给出了制作和评论录屏及其他视频的实用技巧；
\figref{f:online-screencasting}（来自\cite{Chen2009}）复现了那篇论文提出的模式
以及它们之间的关系。
（它也是概念图的一个好例子（\secref{s:memory-concept-maps}）。）

\figpdf{figures/screencast.pdf}{录屏的模式}{f:online-screencasting}{}

那么，是什么让一段教学视频有效？
\cite{Guo2014}通过观察学习者看 MOOC 视频的时长来衡量投入度，
发现：

\begin{itemize}

\item
  更短的视频吸引人得多——视频不应超过六分钟。

\item
  在幻灯片上叠加一个出镜讲解的人，比只有配音配幻灯片更吸引人。

\item
  感觉更个人的视频，可能比高质量录音棚录制更吸引人，
  所以在非正式环境里拍摄，可能比专业棚拍效果更好、成本更低。

\item
  在平板上手写，比 PowerPoint 幻灯片或代码录屏更吸引人，
  不过不清楚这是因为其中的动感和随意，
  还是因为它减少了屏幕上的文字量。

\item
  只要老师热情，讲得相当快也没关系。

\end{itemize}

\cite{Guo2014}没有谈到的一件事，是先有鸡还是先有蛋：
学习者觉得某类视频吸引人，是因为他们习惯了它，
所以生产更多这类视频，会仅仅因为反馈循环而提升投入度吗？
还是说这些建议反映了某种更深的认知过程？
这篇文章也没有考察学习成效：
我们知道学习者对课程的评价与学习并不相关~\cite{Star2014,Uttl2017}，
而且虽然「学习者不会从自己不看的东西里学到东西」说得通，
但他们\emph{确实}会从自己\emph{看}的东西里学到东西，这一点仍有待证明。

\begin{aside}{我有点不安}
  \cite{Guo2014}的研究得到了大学研究伦理委员会的批准，
  那些观看习惯被监测的学习者，几乎肯定在某个时候点过服务条款协议上的「同意」，
  而这些洞见我也乐意知道。
  另一方面，
  在这些结果被发表的会议上，几十篇论文或海报里，
  \emph{没}有任何一篇的标题或摘要里出现过「隐私」一词。
  如果能选，
  我宁愿不知道学习者有多投入，
  也不愿在课堂里助长无处不在的监视。
\end{aside}

录制视频课有很多不同的方式；
为了弄清哪些最有效，
\cite{Mull2007a}把 364 名大学一年级物理学习者
分配到关于牛顿第一、第二定律的在线多媒体处理中，方式有四种：

\begin{description}

\item[讲解：]
  简洁的讲座式呈现。

\item[扩展讲解：]
  在上面基础上加上额外的有趣信息。

\item[反驳：]
  讲解中明确陈述并反驳常见的误解。

\item[对话：]
  学习者—导师围绕与「反驳」相同材料的讨论。

\end{description}

相比「讲解」，「反驳」和「对话」带来最大的学习收益；
先备知识低的学习者受益最大，
而先备知识高的人也没吃亏。
再说一遍，
这凸显了直接处理学习者误解的重要性。
不要只告诉人们\emph{是}什么：
还要告诉他们\emph{不是}什么、以及为什么不是。

\seclbl{混合模式}{s:online-hybrid}
\index{混合式教学}

完全自动化教学只是在教学中使用网络的一种方式。
在实践中，
富裕社会里几乎所有的学习今天都有在线成分，
无论是正式的，
还是通过同伴之间的暗渠道、以及偷偷搜索作业答案。
把现场教学和自动化教学结合起来，能让老师用上两者的长处。
在传统课堂里，
老师能立刻回答问题，
但学习者要等好几天、好几周才能拿到自己编程练习的反馈。
在线的话，
学习者拿到一个答案可能更慢，
但他们能对自己写的代码得到即时反馈
（至少对那些能自动评分的那类练习）。

另一个差别是，
在线练习必须更详细，
因为它们必须预先想到学习者的问题。
我发现，线下课程从「每个人都需要知道的东西」的交集开始，再按需扩展，
而在线课程必须包含「每个人都需要知道的东西」的并集，
因为老师不在场，没人来做扩展。

现实中，
对大多数人来说，在线和线下的区分，
现在已经不如同步和异步的区分重要了：
老师和学习者实时互动吗，
还是他们的交流是分散的、和其他活动交错的？
线下几乎总是同步的，
但在线越来越多地是两者的混合：

\begin{quote}

  我想，我们的孙辈大概会把我们所谓的「真实世界」、
  和他们眼中就是「世界」的那个东西之间的区分，
  看作我们身上最古怪、最不可理解的地方。\\
  —— William Gibson\index{Gibson, William}

\end{quote}

这种混合未来今天最流行的实现，
是\gref{g:flipped-classroom}{翻转课堂}，
也就是学习者自己看录制的课程，
课堂时间用来讨论和做练习。
它最早在~\cite{King1993}里被描述，
作为同伴教学（\secref{s:classroom-peer}）的一部分被推广开来，
并在过去十年里被密集研究。
例如，
\cite{Camp2016}比较了在线学习计算机科学入门课的学习者
和在翻转课堂里学习同一门课的学习者。
（不计分的）练习完成情况，在两组里都和考试分数相关，
但在线学习者的复习练习完成率，
显著低于线下学习者的到课率。

但如果录像可以拿到，
学习者还会来上课做练习吗？
\cite{Nord2017}考察了录像对到课率、
以及不同水平学习者表现的影响。
在大多数情况下，研究发现提供录像没有负面后果；
尤其是，
学习者在有录像可看时并没有翘课
（至少没有比平时翘得更多）。
提供录像的好处对职业生涯早期的学习者最大，
但随着学习者变得更成熟，好处就变小了。

另一种混合模式把在线生活带进课堂。
一起记笔记是第一步（\secref{s:classroom-notetaking}）；
用 \hreffoot{https://www.peardeck.com/}{Pear~Deck}
和 \hreffoot{https://socrative.com/}{Socrative} 之类的工具
实时汇总多项选择题的答案，是另一种。
如果班级很小——比如一打到十五个人——你还可以
让所有学习者加入一个视频会议，
好让他们和老师共享屏幕。
这样他们就能把作品（或问题）展示给全班，
而不必把笔记本电脑接到投影仪上。
学习者随后还能用视频通话里的聊天给老师发问题；
依我的经验，
其中大多数都会由同学来回答，
老师可以等到自然的休息点再处理剩下的。
这种模式有助于为远程学习者拉平起点：
如果有人因为健康原因、
或家庭、工作方面的约束没法到教室，
只要大家都习惯了实时在线协作，
他们仍能几乎平等地参与。

我也用实时远程授课上过课，
学习者们集中在 2—6 个地点、各有助手在场，
而我通过流媒体视频讲课（\secref{s:joining-using}）。
这种模式规模化效果好、
省差旅费，
还能使用结对编程（\secref{s:classroom-pair}）之类的技巧。
\emph{行不通}的，是一组人在现场、一组或多组在远程：
哪怕意愿再好，
现场的参与者也会得到多得多的关注。

\seclbl{在线参与}{s:online-engagement}

\cite{Nuth2007}发现，每间教室里都有三个相互重叠的世界：
公共的（老师在说什么、做什么）、
社会的（学习者之间的同伴互动）、
和私人的（每个学习者的脑子里）。
其中最重要的通常是社会的：
学习者从同伴线索里学到的东西，和从正式教学里学到的一样多。

所以，让任何形式的在线教学有效的关键，
是促进同伴之间的互动。
为帮助做到这一点，
课程几乎总会有某种讨论论坛。
\cite{Mill2016a}观察到，学习者使用它们的方式非常不同：

\begin{quote}

  {\ldots}拖延者特别不可能参与在线讨论论坛，
  而这种更少的参与，
  又和更差的成绩相关。
  对这种相关的一种可能解释是，
  拖延者尤其不愿在讨论已经开始后加入，
  也许因为他们担心在已经成型的对话里被当成新来者。
  这种不愿迟到的加入，
  让他们错失了同伴互动带来的重要学习和动机收益。

\end{quote}

\cite{Vell2017}分析了来自两所大学 395 名 CS2 学生在讨论论坛上的帖子，
把它们分成四类：

\begin{description}

\item[主动：]
  请求帮助，但没有展示推理过程，
  也没有展示学生已经试过什么、或已经知道什么。

\item[建设性：]
  反映学生的推理，
  或尝试构造问题的解。

\item[事务性：]
  课程政策、日程、作业提交等。

\item[内容澄清：]
  请求补充信息，
  但不透露学生自己的想法。

\end{description}

他们发现，建设性和事务性问题占主导，
而且建设性问题与成绩相关。
他们还发现，学生在一门课里很少会问超过一个主动问题，
而这些问题\emph{不}与成绩相关。
虽然这令人失望，
知道这一点有助于设定老师的预期：
我们可能都希望自己的课程有活跃的在线社区，
但我们得接受大多数不会有，
或者大多数学习者之间的讨论会发生在
他们已经在用的、而我们可能并不在其中的渠道里。

\begin{aside}{竞合}
  \cite{Gull2004}描述了一场把协作与竞争结合起来的在线编程比赛。
  比赛在发布一道题目描述、并附上一个正确但低效的解时开始。
  结束时，
  赢家是对改进整个解的总体表现贡献最大的人。
  所有提交都是公开的，
  这样参与者能看到彼此的工作、互相借鉴想法。
  正如论文所示，
  最终的解几乎总是博采众人想法的混血儿。

  \cite{Batt2018}在一门计算入门课里描述了一个小规模的变体。
  第一阶段，
  每个学习者各自提交一个编程项目。
  第二阶段，
  学习者两两结对，为同一个问题造出改进的解。
  评估表明，两阶段项目往往能改善学习者的理解，
  而且他们享受这个过程。
  这类项目不仅提升参与度，
  也让参与者有更多在别人代码基础上继续构建的经验。
\end{aside}

讨论并不是让学习者在线合作的唯一方式。
\cite{Pare2008}和~\cite{Kulk2013}报告了一些实验，
让学习者互相批改作业，
然后把他们给出的分数
与研究生助教或其他专家给的分数作比较。
两项研究都发现，学习者给的分数与专家给的一致，
其一致程度和专家彼此之间的一致程度相当；
而几个简单步骤
（比如过滤掉明显欠考虑的作答、或给评分标准定好结构）
还能进一步减少分歧。
而且正如\secref{s:individual-peer}里讨论的，
串通和偏见在同伴评分里\emph{不是}显著因素。

\begin{aside}{信任，但要教育}
  衡量反馈有效性最常见的办法，
  是把学习者的分数和专家的分数作比较，
  但校准式同伴互评（\secref{s:individual-peer}）可以同样有效。
  \index{校准式同伴互评}
  在让学习者互评作业之前，
  先让他们给样卷打分，并把自己的结果与老师给的分数作比较。
  一旦两者对齐，
  就允许这位学习者开始给同伴打分。
  既然批判性阅读是一种有效的学习方式，
  这个结果或许指向一个未来：
  学习者用技术来做判断，
  而不是被技术评判。
\end{aside}

未来几年我们肯定会更常见到的一种技巧，
是现场编程的在线直播~\cite{Raj2018,Haar2017}。
它具备\secref{s:performance-live}里讨论的大部分好处，
而且当它与协作记笔记（\secref{s:classroom-notetaking}）结合时，
可以很接近课堂内的体验。

再往远处看，
\cite{Ijss2000}识别出在线临场感的四个层次，
从真实感（我们分辨不出差别）
到沉浸（我们忘了有差别）
到投入（我们参与其中，但意识到有差别）
到悬置怀疑（大部分工作由我们自己来完成）。
关键在于，
他们区分了身体临场感，
也就是确实身处某地的感觉，
和社会临场感，也就是与他人同在的感觉。
在大多数学习情境里，后者更重要，
而且再说一遍，
我们可以通过在课堂里使用学习者日常所用的技术来培育它。
例如，
\cite{Deb2018}发现，让学习者用自己的移动设备
对课堂练习做实时反馈，
能改善概念保持和学习者参与度，同时降低不及格率。

在线教学和异步教学都还处在襁褓之中。
集中式的 MOOC 也许会被证明是一条演化的死胡同，
但仍有许多其他有前途的模式值得探索。
尤其是，
\cite{Broo2016}描述了群体进行有成效讨论的五十种方式，
其中只有一小撮被广泛知晓或在线实现。
如果我们顺着学习者的技术习惯走，
而不是要求他们来迁就我们，
我们也许最终学到的东西和他们一样多。

\seclbl{练习}{s:online-exercises}

\exercise{双向视频}{两人一组}{10}

录一段 2—3 分钟的你做某件事的视频，
然后和同伴交换机器，
让你们各自能以 4 倍速看对方的视频。
跟上正在发生的事有多容易？
你有没有漏掉什么？

\exercise{视角}{个人}{10}

根据~\cite{Irib2009}，
不同学科关注影响在线社区成败的不同因素：

\begin{description}

\item[商学：]
  客户忠诚、品牌管理、外在动机。

\item[心理学：]
  社区感、内在动机。

\item[社会学：]
  群体身份、实体社区、社会资本、集体行动。

\item[计算机科学：]
  技术实现。

\end{description}

这些视角里，哪一个最贴近你自己的？
哪一个你最不认同？

\exercise{帮忙还是伤害}{小组}{30}

\hreffoot{https://www.nytimes.com/2018/01/19/business/online-courses-are-harming-the-students-who-need-the-most-help.html}{Susan Dynarski 在\emph{纽约时报}上的文章}
解释了学校如何、以及为何把线下课程不及格的学生塞进在线课程，
以及这如何让他们更容易进一步失败。
读这篇文章，然后：

\begin{enumerate}

\item
  以小组为单位，
  想出 2—3 件学校可以做、以抵消这些负面影响的事，
  并对每人成本做个粗略估计。

\item
  把你的建议和成本与其他组的比较。
  在你看来，对一百名学习者来说，为了腾出资源来落实最受欢迎的那些想法，
  得砍掉多少个全职教学岗位？

\item
  作为一个班，
  你们认为这对学习者是不是净收益？

\end{enumerate}

像这样编预算的练习，是判断谁对教育变革是认真的一个好办法。
人人都能想出他们想做的事；
愿意谈「要促成改变需要哪些取舍」的人，就少得多了。
